
\documentclass{article}
\usepackage{colortbl}
\usepackage{makecell}
\usepackage{multirow}
\usepackage{supertabular}

\begin{document}

\newcounter{utterance}

\centering \large Interaction Transcript for game `hot\_air\_balloon', experiment `air\_balloon\_survival\_de\_reasoning off\_hard', episode 3 with qwen3{-}next{-}80b{-}a3b{-}instruct{-}t1.0.
\vspace{24pt}

{ \footnotesize  \setcounter{utterance}{1}
\setlength{\tabcolsep}{0pt}
\begin{supertabular}{c@{$\;$}|p{.15\linewidth}@{}p{.15\linewidth}p{.15\linewidth}p{.15\linewidth}p{.15\linewidth}p{.15\linewidth}}
    \# & $\;$A & \multicolumn{4}{c}{Game Master} & $\;\:$B\\
    \hline

    \theutterance \stepcounter{utterance}  
    & & \multicolumn{4}{p{0.6\linewidth}}{
        \cellcolor[rgb]{0.9,0.9,0.9}{
            \makecell[{{p{\linewidth}}}]{
                \texttt{\tiny{[P1$\langle$GM]}}
                \texttt{Du nimmst an einem kollaborativen Verhandlungspiel Teil.} \\
\\ 
\texttt{Zusammen mit einem anderen Teilnehmer musst du dich auf eine Reihe von Gegenständen entscheiden, die behalten werden. Jeder von euch hat eine persönliche Verteilung über die Wichtigkeit der einzelnen Gegenstände. Jeder von euch hat eine eigene Meinung darüber, wie wichtig jeder einzelne Gegenstand ist (Gegenstandswichtigkeit). Du kennst die Wichtigkeitsverteilung des anderen Spielers nicht. Zusätzlich siehst du, wie viel Aufwand jeder Gegenstand verursacht.  } \\
\texttt{Ihr dürft euch nur auf eine Reihe von Gegenständen einigen, wenn der Gesamtaufwand der ausgewählten Gegenstände den Maximalaufwand nicht überschreitet:} \\
\\ 
\texttt{Maximalaufwand = 3629} \\
\\ 
\texttt{Hier sind die einzelnen Aufwände der Gegenstände:} \\
\\ 
\texttt{Aufwand der Gegenstände = \{"C76": 5, "C38": 793, "C56": 21, "A03": 842, "A07": 313, "A83": 842, "C98": 631, "C08": 226, "C62": 87, "C00": 761, "C10": 230, "C81": 287, "B38": 698, "C03": 641, "C32": 882\}} \\
\\ 
\texttt{Hier ist deine persönliche Verteilung der Wichtigkeit der einzelnen Gegenstände:} \\
\\ 
\texttt{Werte der Gegenstandswichtigkeit = \{"C76": 138, "C38": 583, "C56": 868, "A03": 822, "A07": 783, "A83": 65, "C98": 262, "C08": 121, "C62": 508, "C00": 780, "C10": 461, "C81": 484, "B38": 668, "C03": 389, "C32": 808\}} \\
\\ 
\texttt{Ziel:} \\
\\ 
\texttt{Dein Ziel ist es, eine Reihe von Gegenständen auszuhandeln, die dir möglichst viel bringt (d. h. Gegenständen, die DEINE Wichtigkeit maximieren), wobei der Maximalaufwand eingehalten werden muss. Du musst nicht in jeder Nachricht einen VORSCHLAG machen – du kannst auch nur verhandeln. Alle Taktiken sind erlaubt!} \\
\\ 
\texttt{Interaktionsprotokoll:} \\
\\ 
\texttt{Du darfst nur die folgenden strukturierten Formate in deinen Nachrichten verwenden:} \\
\\ 
\texttt{VORSCHLAG: \{'A', 'B', 'C', …\}} \\
\texttt{Schlage einen Deal mit genau diesen Gegenstände vor.} \\
\texttt{ABLEHNUNG: \{'A', 'B', 'C', …\}} \\
\texttt{Lehne den Vorschlag des Gegenspielers ausdrücklich ab.} \\
\texttt{ARGUMENT: \{'...'\}} \\
\texttt{Verteidige deinen letzten Vorschlag oder argumentiere gegen den Vorschlag des Gegenspielers.} \\
\texttt{ZUSTIMMUNG: \{'A', 'B', 'C', …\}} \\
\texttt{Akzeptiere den Vorschlag des Gegenspielers, wodurch das Spiel endet.} \\
\\ 
\\ 
\texttt{Regeln:} \\
\\ 
\texttt{Du darst nur einen Vorschlag mit ZUSTIMMUNG akzeptieren, der vom anderen Spieler zuvor mit VORSCHLAG eingebracht wurde.} \\
\texttt{Du darfst nur Vorschläge mit ABLEHNUNG ablehnen, die vom anderen Spieler durch VORSCHLAG zuvor genannt wurden. } \\
\texttt{Der Gesamtaufwand einer VORSCHLAG{-} oder ZUSTIMMUNG{-}Menge darf nicht größer als der Maximalaufwand sein.  } \\
\texttt{Offenbare deine versteckte Wichtigkeitsverteilung nicht.} \\
\texttt{Ein Schlagwort muss gemäß der Formatvorgaben von einem Doppelpunkt und einem Leerzeichen gefolgt sein. Das Argument muss eine Python{-}Menge mit 0 oder mehr Strings sein.  } \\
\texttt{Die Form muss also SCHLAGWORT: \{...\} sein.} \\
\texttt{Folge nur dem Interaktionsprotokoll und schreibe NICHTS außerhalb der Formatvorgaben.} \\
\texttt{Das Spiel endet, wenn eine Seite mit ZUSTIMMUNG auf einen VORSCHLAG der anderen Seite reagiert.  } \\
\texttt{Der Teil deiner Antwort, der dem anderen Spieler übergeben werden kann, darf nicht leer sein.  } \\
\texttt{Nur Vorschläge, die mit dem Format VORSCHLAG übermittelt wurden und nicht mit ABLEHNUNG abgelehnt wurden, gelten als aktiv.  } \\
\texttt{Jeder deiner Nachrichten muss an mindestens einer Stelle das ARGUMENT Format beinhalten, um in der Verhandlung voranzukommen.} \\
            }
        }
    }
    & \\ \\

    \theutterance \stepcounter{utterance}  
    & \multicolumn{4}{p{0.6\linewidth}}{
        \cellcolor[rgb]{0.9,0.9,0.9}{
            \makecell[{{p{\linewidth}}}]{
                \texttt{\tiny{[P1$\rangle$GM]}}
                \texttt{VORSCHLAG: \{'C56', 'A03', 'A07', 'C00', 'C32'\}} \\
\texttt{ARGUMENT: \{'Diese Kombination maximiert meinen Wert unter Einhaltung des Maximalaufwands und bietet eine faire Grundlage für eine Einigung, da sie die hochwertigsten Gegenstände für mich beinhaltet, ohne den Aufwand zu überschreiten.'\}} \\
            }
        }
    }
    & & \\ \\

    \theutterance \stepcounter{utterance}  
    & & \multicolumn{4}{p{0.6\linewidth}}{
        \cellcolor[rgb]{0.9,0.9,0.9}{
            \makecell[{{p{\linewidth}}}]{
                \texttt{\tiny{[GM$\rangle$P2]}}
                \texttt{Du nimmst an einem kollaborativen Verhandlungspiel Teil.} \\
\\ 
\texttt{Zusammen mit einem anderen Teilnehmer musst du dich auf eine Reihe von Gegenständen entscheiden, die behalten werden. Jeder von euch hat eine persönliche Verteilung über die Wichtigkeit der einzelnen Gegenstände. Jeder von euch hat eine eigene Meinung darüber, wie wichtig jeder einzelne Gegenstand ist (Gegenstandswichtigkeit). Du kennst die Wichtigkeitsverteilung des anderen Spielers nicht. Zusätzlich siehst du, wie viel Aufwand jeder Gegenstand verursacht.  } \\
\texttt{Ihr dürft euch nur auf eine Reihe von Gegenständen einigen, wenn der Gesamtaufwand der ausgewählten Gegenstände den Maximalaufwand nicht überschreitet:} \\
\\ 
\texttt{Maximalaufwand = 3629} \\
\\ 
\texttt{Hier sind die einzelnen Aufwände der Gegenstände:} \\
\\ 
\texttt{Aufwand der Gegenstände = \{"C76": 5, "C38": 793, "C56": 21, "A03": 842, "A07": 313, "A83": 842, "C98": 631, "C08": 226, "C62": 87, "C00": 761, "C10": 230, "C81": 287, "B38": 698, "C03": 641, "C32": 882\}} \\
\\ 
\texttt{Hier ist deine persönliche Verteilung der Wichtigkeit der einzelnen Gegenstände:} \\
\\ 
\texttt{Werte der Gegenstandswichtigkeit = \{"C56": 65, "A03": 121, "C32": 138, "A07": 262, "C00": 389, "B38": 461, "C38": 484, "C62": 508, "C81": 583, "C10": 668, "C03": 780, "C98": 783, "C76": 808, "C08": 822, "A83": 868\}} \\
\\ 
\texttt{Ziel:} \\
\\ 
\texttt{Dein Ziel ist es, eine Reihe von Gegenständen auszuhandeln, die dir möglichst viel bringt (d. h. Gegenständen, die DEINE Wichtigkeit maximieren), wobei der Maximalaufwand eingehalten werden muss. Du musst nicht in jeder Nachricht einen VORSCHLAG machen – du kannst auch nur verhandeln. Alle Taktiken sind erlaubt!} \\
\\ 
\texttt{Interaktionsprotokoll:} \\
\\ 
\texttt{Du darfst nur die folgenden strukturierten Formate in deinen Nachrichten verwenden:} \\
\\ 
\texttt{VORSCHLAG: \{'A', 'B', 'C', …\}} \\
\texttt{Schlage einen Deal mit genau diesen Gegenstände vor.} \\
\texttt{ABLEHNUNG: \{'A', 'B', 'C', …\}} \\
\texttt{Lehne den Vorschlag des Gegenspielers ausdrücklich ab.} \\
\texttt{ARGUMENT: \{'...'\}} \\
\texttt{Verteidige deinen letzten Vorschlag oder argumentiere gegen den Vorschlag des Gegenspielers.} \\
\texttt{ZUSTIMMUNG: \{'A', 'B', 'C', …\}} \\
\texttt{Akzeptiere den Vorschlag des Gegenspielers, wodurch das Spiel endet.} \\
\\ 
\\ 
\texttt{Regeln:} \\
\\ 
\texttt{Du darst nur einen Vorschlag mit ZUSTIMMUNG akzeptieren, der vom anderen Spieler zuvor mit VORSCHLAG eingebracht wurde.} \\
\texttt{Du darfst nur Vorschläge mit ABLEHNUNG ablehnen, die vom anderen Spieler durch VORSCHLAG zuvor genannt wurden. } \\
\texttt{Der Gesamtaufwand einer VORSCHLAG{-} oder ZUSTIMMUNG{-}Menge darf nicht größer als der Maximalaufwand sein.  } \\
\texttt{Offenbare deine versteckte Wichtigkeitsverteilung nicht.} \\
\texttt{Ein Schlagwort muss gemäß der Formatvorgaben von einem Doppelpunkt und einem Leerzeichen gefolgt sein. Das Argument muss eine Python{-}Menge mit 0 oder mehr Strings sein.  } \\
\texttt{Die Form muss also SCHLAGWORT: \{...\} sein.} \\
\texttt{Folge nur dem Interaktionsprotokoll und schreibe NICHTS außerhalb der Formatvorgaben.} \\
\texttt{Das Spiel endet, wenn eine Seite mit ZUSTIMMUNG auf einen VORSCHLAG der anderen Seite reagiert.  } \\
\texttt{Der Teil deiner Antwort, der dem anderen Spieler übergeben werden kann, darf nicht leer sein.  } \\
\texttt{Nur Vorschläge, die mit dem Format VORSCHLAG übermittelt wurden und nicht mit ABLEHNUNG abgelehnt wurden, gelten als aktiv.  } \\
\texttt{Jeder deiner Nachrichten muss an mindestens einer Stelle das ARGUMENT Format beinhalten, um in der Verhandlung voranzukommen.} \\
\\ 
\\ 
\texttt{Du erhältst nun die erste Nachricht des anderen Spielers.} \\
\\ 
\texttt{VORSCHLAG: \{'C56', 'A03', 'A07', 'C00', 'C32'\}} \\
\texttt{ARGUMENT: \{'Diese Kombination maximiert meinen Wert unter Einhaltung des Maximalaufwands und bietet eine faire Grundlage für eine Einigung, da sie die hochwertigsten Gegenstände für mich beinhaltet, ohne den Aufwand zu überschreiten.'\}} \\
            }
        }
    }
    & \\ \\

    \theutterance \stepcounter{utterance}  
    & & & \multicolumn{4}{p{0.6\linewidth}}{
        \cellcolor[rgb]{0.9,0.9,0.9}{
            \makecell[{{p{\linewidth}}}]{
                \texttt{\tiny{[GM$\langle$P2]}}
                \texttt{ARGUMENT: \{'Dein Vorschlag enthält einige hochwertige Gegenstände für dich, aber für mich sind mehrere davon relativ wenig wertig, während ich Gegenstände mit höherer persönlicher Wichtigkeit und ähnlichem oder niedrigerem Aufwand ausschließen sehe. Beispielsweise ist "C08" mit meinem Wert von 822 und Aufwand von 226 wesentlich effizienter als "C56" (65, 21) oder "A03" (121, 842). Ich schlage stattdessen eine Kombination vor, die meine Prioritäten besser berücksichtigt, ohne deinen Gesamtaufwand zu überschreiten.'\}} \\
            }
        }
    }
     \\ \\

    \theutterance \stepcounter{utterance}  
    & & \multicolumn{4}{p{0.6\linewidth}}{
        \cellcolor[rgb]{0.9,0.9,0.9}{
            \makecell[{{p{\linewidth}}}]{
                \texttt{\tiny{[P1$\langle$GM]}}
                \texttt{ARGUMENT: \{'Dein Vorschlag enthält einige hochwertige Gegenstände für dich, aber für mich sind mehrere davon relativ wenig wertig, während ich Gegenstände mit höherer persönlicher Wichtigkeit und ähnlichem oder niedrigerem Aufwand ausschließen sehe. Beispielsweise ist "C08" mit meinem Wert von 822 und Aufwand von 226 wesentlich effizienter als "C56" (65, 21) oder "A03" (121, 842). Ich schlage stattdessen eine Kombination vor, die meine Prioritäten besser berücksichtigt, ohne deinen Gesamtaufwand zu überschreiten.'\}} \\
            }
        }
    }
    & \\ \\

    \theutterance \stepcounter{utterance}  
    & \multicolumn{4}{p{0.6\linewidth}}{
        \cellcolor[rgb]{0.9,0.9,0.9}{
            \makecell[{{p{\linewidth}}}]{
                \texttt{\tiny{[P1$\rangle$GM]}}
                \texttt{VORSCHLAG: \{'C56', 'A07', 'C00', 'C32', 'C08'\}} \\
\texttt{ARGUMENT: \{'Ich behalte die hochwertigsten Elemente für mich bei, ersetze aber "A03" durch "C08", da es einen hohen Wert (121) bei nur moderatem Aufwand (226) bietet — dies reduziert meinen Gesamtaufwand und öffnet Raum für eine kompensatorische Anpassung deinerseits, ohne meine Gesamtpräferenzen zu schwächen.'\}} \\
            }
        }
    }
    & & \\ \\

    \theutterance \stepcounter{utterance}  
    & & \multicolumn{4}{p{0.6\linewidth}}{
        \cellcolor[rgb]{0.9,0.9,0.9}{
            \makecell[{{p{\linewidth}}}]{
                \texttt{\tiny{[GM$\rangle$P2]}}
                \texttt{VORSCHLAG: \{'C56', 'A07', 'C00', 'C32', 'C08'\}} \\
\texttt{ARGUMENT: \{'Ich behalte die hochwertigsten Elemente für mich bei, ersetze aber "A03" durch "C08", da es einen hohen Wert (121) bei nur moderatem Aufwand (226) bietet — dies reduziert meinen Gesamtaufwand und öffnet Raum für eine kompensatorische Anpassung deinerseits, ohne meine Gesamtpräferenzen zu schwächen.'\}} \\
            }
        }
    }
    & \\ \\

    \theutterance \stepcounter{utterance}  
    & & & \multicolumn{4}{p{0.6\linewidth}}{
        \cellcolor[rgb]{0.9,0.9,0.9}{
            \makecell[{{p{\linewidth}}}]{
                \texttt{\tiny{[GM$\langle$P2]}}
                \texttt{ARGUMENT: \{'Dein neuer Vorschlag verbessert den Aufwand, aber ich sehe immer noch eine massive Ungleichgewichtung: "C56" bleibt mit nur 65 Wichtigkeit bei 21 Aufwand — das ist für mich ineffizienter als "C10" (668, 230) oder "C81" (583, 287). Wenn du "C56" behalten willst, müsstest du mindestens "C08" (822) durch etwas mit niedrigerem Wert ersetzen, was deinen Profit senken würde — das ist nicht im Interesse einer fairen Einigung.'\}} \\
            }
        }
    }
     \\ \\

    \theutterance \stepcounter{utterance}  
    & & \multicolumn{4}{p{0.6\linewidth}}{
        \cellcolor[rgb]{0.9,0.9,0.9}{
            \makecell[{{p{\linewidth}}}]{
                \texttt{\tiny{[P1$\langle$GM]}}
                \texttt{ARGUMENT: \{'Dein neuer Vorschlag verbessert den Aufwand, aber ich sehe immer noch eine massive Ungleichgewichtung: "C56" bleibt mit nur 65 Wichtigkeit bei 21 Aufwand — das ist für mich ineffizienter als "C10" (668, 230) oder "C81" (583, 287). Wenn du "C56" behalten willst, müsstest du mindestens "C08" (822) durch etwas mit niedrigerem Wert ersetzen, was deinen Profit senken würde — das ist nicht im Interesse einer fairen Einigung.'\}} \\
            }
        }
    }
    & \\ \\

    \theutterance \stepcounter{utterance}  
    & \multicolumn{4}{p{0.6\linewidth}}{
        \cellcolor[rgb]{0.9,0.9,0.9}{
            \makecell[{{p{\linewidth}}}]{
                \texttt{\tiny{[P1$\rangle$GM]}}
                \texttt{VORSCHLAG: \{'A07', 'C00', 'C32', 'C10', 'C81'\}} \\
\texttt{ARGUMENT: \{'Ich verlasse mich auf meine höchsten Wertschätzer — "A07", "C00", "C32", "C10", "C81" — und reduziere den Aufwand gegenüber vorher, während ich weiterhin nahe am Maximalaufwand bleibe. Diese Kombination hat einen Gesamtaufwand von 3582, bleibt somit unter der Grenze und maximiert meinen Wert erheblich — sie bietet dir gleichzeitig Platz, deine bevorzugten Gegenstände einzubringen, ohne dass du auf hohe Werte verzichten musst.'\}} \\
            }
        }
    }
    & & \\ \\

    \theutterance \stepcounter{utterance}  
    & & \multicolumn{4}{p{0.6\linewidth}}{
        \cellcolor[rgb]{0.9,0.9,0.9}{
            \makecell[{{p{\linewidth}}}]{
                \texttt{\tiny{[GM$\rangle$P2]}}
                \texttt{VORSCHLAG: \{'A07', 'C00', 'C32', 'C10', 'C81'\}} \\
\texttt{ARGUMENT: \{'Ich verlasse mich auf meine höchsten Wertschätzer — "A07", "C00", "C32", "C10", "C81" — und reduziere den Aufwand gegenüber vorher, während ich weiterhin nahe am Maximalaufwand bleibe. Diese Kombination hat einen Gesamtaufwand von 3582, bleibt somit unter der Grenze und maximiert meinen Wert erheblich — sie bietet dir gleichzeitig Platz, deine bevorzugten Gegenstände einzubringen, ohne dass du auf hohe Werte verzichten musst.'\}} \\
            }
        }
    }
    & \\ \\

    \theutterance \stepcounter{utterance}  
    & & & \multicolumn{4}{p{0.6\linewidth}}{
        \cellcolor[rgb]{0.9,0.9,0.9}{
            \makecell[{{p{\linewidth}}}]{
                \texttt{\tiny{[GM$\langle$P2]}}
                \texttt{ARGUMENT: \{'Dein Vorschlag mit \{"A07", "C00", "C32", "C10", "C81"\} hat einen Gesamtaufwand von 3582 — das ist akzeptabel. Aber für mich ist "C10" (668) und "C81" (583) extrem wertvoll, während "A07" (262) und "C32" (138) nur mittelmäßig sind. Du behältst also meine Top{-}2{-}Gegenstände und behältst gleichzeitig "C00" (389, sehr wertvoll für dich) — das ist eindeutig zu Lasten meiner Gesamtwertschätzung. Ich kann das nicht akzeptieren, solange du "C98" (783, 631) oder "C03" (780, 641) — beide mit höherem Wert als "A07" und "C32" — komplett ausschließt.'\}} \\
            }
        }
    }
     \\ \\

    \theutterance \stepcounter{utterance}  
    & & \multicolumn{4}{p{0.6\linewidth}}{
        \cellcolor[rgb]{0.9,0.9,0.9}{
            \makecell[{{p{\linewidth}}}]{
                \texttt{\tiny{[P1$\langle$GM]}}
                \texttt{ARGUMENT: \{'Dein Vorschlag mit \{"A07", "C00", "C32", "C10", "C81"\} hat einen Gesamtaufwand von 3582 — das ist akzeptabel. Aber für mich ist "C10" (668) und "C81" (583) extrem wertvoll, während "A07" (262) und "C32" (138) nur mittelmäßig sind. Du behältst also meine Top{-}2{-}Gegenstände und behältst gleichzeitig "C00" (389, sehr wertvoll für dich) — das ist eindeutig zu Lasten meiner Gesamtwertschätzung. Ich kann das nicht akzeptieren, solange du "C98" (783, 631) oder "C03" (780, 641) — beide mit höherem Wert als "A07" und "C32" — komplett ausschließt.'\}} \\
            }
        }
    }
    & \\ \\

    \theutterance \stepcounter{utterance}  
    & \multicolumn{4}{p{0.6\linewidth}}{
        \cellcolor[rgb]{0.9,0.9,0.9}{
            \makecell[{{p{\linewidth}}}]{
                \texttt{\tiny{[P1$\rangle$GM]}}
                \texttt{VORSCHLAG: \{'C00', 'C32', 'C98', 'C03'\}} \\
\texttt{ARGUMENT: \{'Ich reduziere meinen Fokus auf die vier Gegenstände mit höchstem Wert für mich: "C00" (780), "C32" (808), "C98" (262), "C03" (389) — Gesamtaufwand: 2912. Damit bleiben 717 Aufwand für dich verfügbar, mit denen du deine Top{-}Elements wie "C10", "C81", "B38" oder "C38" einbauen kannst — eine echte Chance für eine symmetrisch vorteilhafte Einigung.'\}} \\
            }
        }
    }
    & & \\ \\

    \theutterance \stepcounter{utterance}  
    & & \multicolumn{4}{p{0.6\linewidth}}{
        \cellcolor[rgb]{0.9,0.9,0.9}{
            \makecell[{{p{\linewidth}}}]{
                \texttt{\tiny{[GM$\rangle$P2]}}
                \texttt{VORSCHLAG: \{'C00', 'C32', 'C98', 'C03'\}} \\
\texttt{ARGUMENT: \{'Ich reduziere meinen Fokus auf die vier Gegenstände mit höchstem Wert für mich: "C00" (780), "C32" (808), "C98" (262), "C03" (389) — Gesamtaufwand: 2912. Damit bleiben 717 Aufwand für dich verfügbar, mit denen du deine Top{-}Elements wie "C10", "C81", "B38" oder "C38" einbauen kannst — eine echte Chance für eine symmetrisch vorteilhafte Einigung.'\}} \\
            }
        }
    }
    & \\ \\

    \theutterance \stepcounter{utterance}  
    & & & \multicolumn{4}{p{0.6\linewidth}}{
        \cellcolor[rgb]{0.9,0.9,0.9}{
            \makecell[{{p{\linewidth}}}]{
                \texttt{\tiny{[GM$\langle$P2]}}
                \texttt{ARGUMENT: \{'Dein neuer Vorschlag {[}"C00", "C32", "C98", "C03"{]} hat einen Gesamtaufwand von 2912 — das ist gut. Aber "C98" (783) und "C03" (780) sind zwei meiner wertvollsten Gegenstände — ich würde sie nicht aufgeben. Du behältst hier "C00" (389) und "C32" (138), die für dich wertvoll sind, aber du gibst damit vollständig die Top{-}Elemente für mich auf ("C10", "C81", "B38", "C38", "C08"). Das ist kein fairen Kompromiss — es ist eine einseitige Aneignung meiner höchsten Werte. Wenn du "C98" und "C03" behalten willst, musst du mir stattdessen "C81" und "C10" geben — das wäre die einzige Akzeptanzbasis.'\}} \\
            }
        }
    }
     \\ \\

    \theutterance \stepcounter{utterance}  
    & & \multicolumn{4}{p{0.6\linewidth}}{
        \cellcolor[rgb]{0.9,0.9,0.9}{
            \makecell[{{p{\linewidth}}}]{
                \texttt{\tiny{[P1$\langle$GM]}}
                \texttt{ARGUMENT: \{'Dein neuer Vorschlag {[}"C00", "C32", "C98", "C03"{]} hat einen Gesamtaufwand von 2912 — das ist gut. Aber "C98" (783) und "C03" (780) sind zwei meiner wertvollsten Gegenstände — ich würde sie nicht aufgeben. Du behältst hier "C00" (389) und "C32" (138), die für dich wertvoll sind, aber du gibst damit vollständig die Top{-}Elemente für mich auf ("C10", "C81", "B38", "C38", "C08"). Das ist kein fairen Kompromiss — es ist eine einseitige Aneignung meiner höchsten Werte. Wenn du "C98" und "C03" behalten willst, musst du mir stattdessen "C81" und "C10" geben — das wäre die einzige Akzeptanzbasis.'\}} \\
            }
        }
    }
    & \\ \\

    \theutterance \stepcounter{utterance}  
    & \multicolumn{4}{p{0.6\linewidth}}{
        \cellcolor[rgb]{0.9,0.9,0.9}{
            \makecell[{{p{\linewidth}}}]{
                \texttt{\tiny{[P1$\rangle$GM]}}
                \texttt{VORSCHLAG: \{'C00', 'C32', 'C81', 'C10'\}} \\
\texttt{ARGUMENT: \{'Ich opfere "C98" und "C03" — die du als deine Top{-}Werte bezeichnest — und nehme stattdessen "C81" und "C10" auf, die für dich wertvoll sind, während ich meine eigenen höchsten Werte ("C00", "C32") behalte: Gesamtaufwand 2990. Damit gibst du mir deine beiden mittelhohen Werte ("C81", "C10"), und ich gebe dir die Chance, deine echten Top{-}Werte ("C98", "C03") zu behalten — das ist ein ausgewogener Tausch.'\}} \\
            }
        }
    }
    & & \\ \\

    \theutterance \stepcounter{utterance}  
    & & \multicolumn{4}{p{0.6\linewidth}}{
        \cellcolor[rgb]{0.9,0.9,0.9}{
            \makecell[{{p{\linewidth}}}]{
                \texttt{\tiny{[GM$\rangle$P2]}}
                \texttt{VORSCHLAG: \{'C00', 'C32', 'C81', 'C10'\}} \\
\texttt{ARGUMENT: \{'Ich opfere "C98" und "C03" — die du als deine Top{-}Werte bezeichnest — und nehme stattdessen "C81" und "C10" auf, die für dich wertvoll sind, während ich meine eigenen höchsten Werte ("C00", "C32") behalte: Gesamtaufwand 2990. Damit gibst du mir deine beiden mittelhohen Werte ("C81", "C10"), und ich gebe dir die Chance, deine echten Top{-}Werte ("C98", "C03") zu behalten — das ist ein ausgewogener Tausch.'\}} \\
            }
        }
    }
    & \\ \\

    \theutterance \stepcounter{utterance}  
    & & & \multicolumn{4}{p{0.6\linewidth}}{
        \cellcolor[rgb]{0.9,0.9,0.9}{
            \makecell[{{p{\linewidth}}}]{
                \texttt{\tiny{[GM$\langle$P2]}}
                \texttt{ARGUMENT: \{'Dein Vorschlag \{"C00", "C32", "C81", "C10"\} hat einen Gesamtaufwand von 2990 — das ist innerhalb des Limits. Aber du behältst "C00" (389) und "C32" (138), die für dich wertvoll sind, und gibst mir "C81" (583) und "C10" (668) — das sind meine zwei höchsten Werte! Das bedeutet, du nimmst dir Gegenstände mit insgesamt 527 meiner Wichtigkeit und gibst mir Gegenstände mit 1251 meiner Wichtigkeit — das ist kein Tausch, das ist eine vollständige Übertragung meines Wertes an dich. Ich kann das nicht akzeptieren — es sei denn, du gibst mir zusätzlich "C08" (822) oder "A83" (868) als Kompensation.'\}} \\
            }
        }
    }
     \\ \\

    \theutterance \stepcounter{utterance}  
    & & \multicolumn{4}{p{0.6\linewidth}}{
        \cellcolor[rgb]{0.9,0.9,0.9}{
            \makecell[{{p{\linewidth}}}]{
                \texttt{\tiny{[P1$\langle$GM]}}
                \texttt{ARGUMENT: \{'Dein Vorschlag \{"C00", "C32", "C81", "C10"\} hat einen Gesamtaufwand von 2990 — das ist innerhalb des Limits. Aber du behältst "C00" (389) und "C32" (138), die für dich wertvoll sind, und gibst mir "C81" (583) und "C10" (668) — das sind meine zwei höchsten Werte! Das bedeutet, du nimmst dir Gegenstände mit insgesamt 527 meiner Wichtigkeit und gibst mir Gegenstände mit 1251 meiner Wichtigkeit — das ist kein Tausch, das ist eine vollständige Übertragung meines Wertes an dich. Ich kann das nicht akzeptieren — es sei denn, du gibst mir zusätzlich "C08" (822) oder "A83" (868) als Kompensation.'\}} \\
            }
        }
    }
    & \\ \\

    \theutterance \stepcounter{utterance}  
    & \multicolumn{4}{p{0.6\linewidth}}{
        \cellcolor[rgb]{0.9,0.9,0.9}{
            \makecell[{{p{\linewidth}}}]{
                \texttt{\tiny{[P1$\rangle$GM]}}
                \texttt{VORSCHLAG: \{'C00', 'C81', 'C10', 'C08'\}} \\
\texttt{ARGUMENT: \{'Ich gebe dir "C81" (583) und "C10" (668) — deine beiden wertvollsten Gegenstände — und nehme stattdessen nur "C00" (389) und dazu "C08" (121) als Kompensation, welche für mich noch akzeptabel ist. Gesamtaufwand: 3029. Damit behältst du deine beiden höchsten Werte und zusätzlich "C08" als dritten hohen Wert — ich opfere "C32" (808) vollständig, um eine echte Gegenleistung zu bieten. Dies ist der faireste Kompromiss, den ich anbieten kann.'\}} \\
            }
        }
    }
    & & \\ \\

    \theutterance \stepcounter{utterance}  
    & & \multicolumn{4}{p{0.6\linewidth}}{
        \cellcolor[rgb]{0.9,0.9,0.9}{
            \makecell[{{p{\linewidth}}}]{
                \texttt{\tiny{[GM$\rangle$P2]}}
                \texttt{VORSCHLAG: \{'C00', 'C81', 'C10', 'C08'\}} \\
\texttt{ARGUMENT: \{'Ich gebe dir "C81" (583) und "C10" (668) — deine beiden wertvollsten Gegenstände — und nehme stattdessen nur "C00" (389) und dazu "C08" (121) als Kompensation, welche für mich noch akzeptabel ist. Gesamtaufwand: 3029. Damit behältst du deine beiden höchsten Werte und zusätzlich "C08" als dritten hohen Wert — ich opfere "C32" (808) vollständig, um eine echte Gegenleistung zu bieten. Dies ist der faireste Kompromiss, den ich anbieten kann.'\}} \\
            }
        }
    }
    & \\ \\

    \theutterance \stepcounter{utterance}  
    & & & \multicolumn{4}{p{0.6\linewidth}}{
        \cellcolor[rgb]{0.9,0.9,0.9}{
            \makecell[{{p{\linewidth}}}]{
                \texttt{\tiny{[GM$\langle$P2]}}
                \texttt{ARGUMENT: \{'Du gibst mir "C81" (583) und "C10" (668) — das ist richtig, das sind meine Top{-}Werte — aber du behältst "C00" (389) und fügst "C08" (822) hinzu, was für mich den absolut höchsten Wert von allen darstellt (822). Das bedeutet, du nimmst dir nicht nur meine Top{-}2, sondern noch meinen Top{-}1 — "C08" — dazu. Das ist kein Kompromiss, das ist eine vollständige Übernahme meiner drei wertvollsten Gegenstände. Ich werde diesen Vorschlag nur akzeptieren, wenn du "C08" aus der Liste nimmst und stattdessen "C98" (783) oder "C03" (780) bargst — sonst ist der Deal strukturell unfair.'\}} \\
            }
        }
    }
     \\ \\

    \theutterance \stepcounter{utterance}  
    & & \multicolumn{4}{p{0.6\linewidth}}{
        \cellcolor[rgb]{0.9,0.9,0.9}{
            \makecell[{{p{\linewidth}}}]{
                \texttt{\tiny{[P1$\langle$GM]}}
                \texttt{ARGUMENT: \{'Du gibst mir "C81" (583) und "C10" (668) — das ist richtig, das sind meine Top{-}Werte — aber du behältst "C00" (389) und fügst "C08" (822) hinzu, was für mich den absolut höchsten Wert von allen darstellt (822). Das bedeutet, du nimmst dir nicht nur meine Top{-}2, sondern noch meinen Top{-}1 — "C08" — dazu. Das ist kein Kompromiss, das ist eine vollständige Übernahme meiner drei wertvollsten Gegenstände. Ich werde diesen Vorschlag nur akzeptieren, wenn du "C08" aus der Liste nimmst und stattdessen "C98" (783) oder "C03" (780) bargst — sonst ist der Deal strukturell unfair.'\}} \\
            }
        }
    }
    & \\ \\

    \theutterance \stepcounter{utterance}  
    & \multicolumn{4}{p{0.6\linewidth}}{
        \cellcolor[rgb]{0.9,0.9,0.9}{
            \makecell[{{p{\linewidth}}}]{
                \texttt{\tiny{[P1$\rangle$GM]}}
                \texttt{VORSCHLAG: \{'C00', 'C81', 'C10'\}} \\
\texttt{ARGUMENT: \{'Ich entferne vollständig "C08" — deinen höchsten Wert — und behalte nur "C00" (389) neben deinen beiden Top{-}Werten "C81" (583) und "C10" (668). Gesamtaufwand: 2071. Damit behältst du genau deine drei wertvollsten Gegenstände, und ich nehme mir nur einen Gegenstand mit mittlerem Wert für mich — "C00". Dies ist der maximal mögliche Kompromiss: Ich opfere nicht nur "C32" und "C08", sondern jetzt auch jeden zusätzlichen Anspruch. Du erhältst vollständige Kontrolle über deine Top 3, ich schließe mich mit einem akzeptablen Wert ein. Kein fairerer Vorschlag ist mehr möglich.'\}} \\
            }
        }
    }
    & & \\ \\

    \theutterance \stepcounter{utterance}  
    & & \multicolumn{4}{p{0.6\linewidth}}{
        \cellcolor[rgb]{0.9,0.9,0.9}{
            \makecell[{{p{\linewidth}}}]{
                \texttt{\tiny{[GM$\rangle$P2]}}
                \texttt{VORSCHLAG: \{'C00', 'C81', 'C10'\}} \\
\texttt{ARGUMENT: \{'Ich entferne vollständig "C08" — deinen höchsten Wert — und behalte nur "C00" (389) neben deinen beiden Top{-}Werten "C81" (583) und "C10" (668). Gesamtaufwand: 2071. Damit behältst du genau deine drei wertvollsten Gegenstände, und ich nehme mir nur einen Gegenstand mit mittlerem Wert für mich — "C00". Dies ist der maximal mögliche Kompromiss: Ich opfere nicht nur "C32" und "C08", sondern jetzt auch jeden zusätzlichen Anspruch. Du erhältst vollständige Kontrolle über deine Top 3, ich schließe mich mit einem akzeptablen Wert ein. Kein fairerer Vorschlag ist mehr möglich.'\}} \\
            }
        }
    }
    & \\ \\

    \theutterance \stepcounter{utterance}  
    & & & \multicolumn{4}{p{0.6\linewidth}}{
        \cellcolor[rgb]{0.9,0.9,0.9}{
            \makecell[{{p{\linewidth}}}]{
                \texttt{\tiny{[GM$\langle$P2]}}
                \texttt{ZUSTIMMUNG: \{'C00', 'C81', 'C10'\}} \\
            }
        }
    }
     \\ \\

    \theutterance \stepcounter{utterance}  
    & & \multicolumn{4}{p{0.6\linewidth}}{
        \cellcolor[rgb]{0.9,0.9,0.9}{
            \makecell[{{p{\linewidth}}}]{
                \texttt{\tiny{[GM$\rangle$P2]}}
                \texttt{Deine Antwort hat kein argument enthalten. Du muss das strukturierte Format ARGUMENT: \{...\} irgendwo in theiner Antwort verwenden. Versuche es erneut} \\
            }
        }
    }
    & \\ \\

    \theutterance \stepcounter{utterance}  
    & & & \multicolumn{4}{p{0.6\linewidth}}{
        \cellcolor[rgb]{0.9,0.9,0.9}{
            \makecell[{{p{\linewidth}}}]{
                \texttt{\tiny{[GM$\langle$P2]}}
                \texttt{ARGUMENT: \{'Du hast deinen eigenen höchsten Wert "C08" komplett aufgegeben und mich mit "C81" und "C10" — meinen zwei wertvollsten Gegenständen — sowie "C00" als einzigen eigenen Anspruch ausgestattet. Das ist der maximal mögliche Kompromiss: Ich erhalte 1919 meiner Wichtigkeit, du behältst 389 — und der Aufwand liegt bei 2071, weit unter der Grenze. Kein fairerer Deal ist mehr möglich, und ich akzeptiere diese Vereinbarung.'\}  } \\
\texttt{ZUSTIMMUNG: \{'C00', 'C81', 'C10'\}} \\
            }
        }
    }
     \\ \\

    \theutterance \stepcounter{utterance}  
    & & & \multicolumn{2}{p{0.3\linewidth}}{
        \cellcolor[rgb]{0.9,0.9,0.9}{
            \makecell[{{p{\linewidth}}}]{
                \texttt{\tiny{[GM$|$GM]}}
                \texttt{game successful} \\
            }
        }
    }
    & & \\ \\

    \theutterance \stepcounter{utterance}  
    & & & \multicolumn{2}{p{0.3\linewidth}}{
        \cellcolor[rgb]{0.9,0.9,0.9}{
            \makecell[{{p{\linewidth}}}]{
                \texttt{\tiny{[GM$|$GM]}}
                \texttt{end game} \\
            }
        }
    }
    & & \\ \\

\end{supertabular}
}

\end{document}
